\subsection{Three Comings}\label{sec:three-comings}

St Bernard began a sermon to his monks on the first day of the season by
saying that its name was known to all the world, ``But perhaps the meaning of
the name is not so generally understood.''

\latin{Adventus} is Latin for ``a coming''. The Church keeps four weeks in
expectation of a coming, and she does it every year, though the Son of God has
already come. St Cyril of Jerusalem, in one of his catechetical lectures, sets
the two comings side by side:

\begin{quote}
For all things, for the most part, are twofold in our Lord Jesus Christ: a
twofold generation; one, of God, before the ages; and one, of a Virgin, at the
close of the ages: His descents twofold; one, the unobserved, like rain on a
fleece; and a second His open coming, which is to be. \dots\ In His first
coming, He endured the Cross, despising shame; in His second, He comes attended
by a host of Angels, receiving glory. We rest not then upon His first advent
only, but look also for His second.\par
{\smallskip\noindent\small ---St Cyril of Jerusalem, Catechetical Lecture 15\par}
\end{quote}

The first coming was hidden, like rain falling on a fleece; the second will be
open to every eye. Advent keeps both, and on the Vigil, at the end of the
season, the Church asks for both in the Collect of the day:

\begin{studyframe}
\latin{Deus, qui nos redemptionis nostrae annua exspectatione laetificas:
praesta; ut Unigenitum tuum, quem Redemptorem laeti suscipimus, venientem
quoque judicem securi videamus, Dominum nostrum Jesum Christum, Filium tuum:
Qui tecum vivit et regnat in unitate.}
\medskip

O God, Who dost gladden us year by year with the expectation of our
redemption, grant that we, who now with joy receive Thine only begotten Son as
our Redeemer, may behold Him also without fear, when He cometh as our judge,
our Lord Jesus Christ. Who with thee.\par
\end{studyframe}

So the First Sunday's Gospel is the Lord's own prophecy of the end: signs in
the sun and moon and stars, the sea roaring, men withering away for fear, and
the Son of man coming in a cloud. It does not end in terror. In the midst of it
comes a word to the disciples:

\begin{quote}
But when these things begin to come to pass, look up and lift up your heads,
because your redemption is at hand.\par
{\smallskip\noindent\small ---Luke 21:28\par}
\end{quote}

St Gregory the Great says what lifting the head means: ``To lift up our heads
then, is to raise up our minds to the joys of the heavenly country.'' So the
Church sets the end of the world on the threshold of Christmas, and bids her
children meet it with their heads up.

\subsubsection*{The hidden coming}

Between those two comings St Bernard found a third. In one of his Advent sermons
he speaks of a coming that is neither the first nor the last, and that only
those who receive it know:

\begin{quote}
This intermediate place is in fact another advent wherein such as have
knowledge thereof enjoy a most pleasant repose. For whereas the other two are
known to all, this is secret. In the first ``He (Christ) was seen on earth and
conversed with men,'' when, as He Himself bears witness, ``they both saw and
hated'' Him. In the last ``all flesh shall see the Salvation of God,'' and ``they
shall look on Him Whom they pierced.'' But in that which intervenes He is hidden,
visible only to the elect who see Him in themselves, and so their souls are
saved.\par
{\smallskip\noindent\small ---St Bernard, Advent Sermon V\par}
\end{quote}

A few lines further on, Bernard calls this middle coming \latin{via quædam},
``a kind of road'', by which one comes from the first to the last: \latin{In
primo Christus fuit redemptio nostra: in ultimo apparebit vita nostra}, ``In the
first, Christ was our redemption; in the last, he will appear as our life''; and
in this middle coming, \latin{requies est et consolatio nostra}, ``he is our
rest and our consolation''. John Beleth, who wrote on the Church's offices,
counted the same three: Christ comes \latin{in carne}, in the flesh;
\latin{in mente}, in the mind, every day, by the Holy Spirit; and \latin{in
majestate}, in majesty, on the day of judgment.

How is a hidden coming received? Bernard's image is homely. The Lord comes
into the soul by his word, and the word is kept the way bread is kept:

\begin{quote}
Material bread, whilst it is kept in the cupboard may be stolen by thieves, or
nibbled by mice, or may become stale and unfit for use. But once it has been
assimilated, we have nothing further to fear from such possibilities. It is thus
we should keep the word of God, for ``blessed are they who keep it'' so. Receive
this spiritual bread, therefore, into the stomach of your minds, and having been
there properly digested, let it be transmitted thence to the affections and the
will.\par
{\smallskip\noindent\small ---St Bernard, Advent Sermon V\par}
\end{quote}

Here is the work of the season for a Christian in the world. The first coming
is past, and the last is in God's hands. The middle one can be refused, because
it comes by a word that may be heard and forgotten, or heard and eaten. Advent
asks for a little of the word each day, taken in and kept: a lesson of Isaiah,
a few lines of a Father, a Collect said slowly in the morning.

The Church did not always keep four weeks for this. She came to them slowly,
by two roads, and she learned on the way to keep them as a fast and as a joy.
